\title{Parameter-Efficient Orthogonal Finetuning via Factorization}

\begin{abstract}
Large foundation models have become prevalent, yet the cost associated with training them from the ground up is exceedingly high. Consequently, finding ways to adapt these sophisticated models to specific downstream applications in an efficient manner is essential. In this study, we explore a structured approach to model adaptation—Orthogonal Finetuning (OFT)—for downstream finetuning tasks. Although OFT offers promising generalization, its reliance on high-dimensional orthogonal matrices means it still involves a substantial number of trainable parameters. To tackle this limitation, we first reinterpret OFT from the angle of information transfer, elucidating several central criteria for enhancing parameter efficiency. Drawing inspiration from the efficiency gains achieved in the Cooley-Tukey fast Fourier transform, we develop a new orthogonal parameterization that leverages butterfly structures. By integrating this parameterization into OFT, we introduce Orthogonal Butterfly (BOFT), a novel approach that prioritizes parameter efficiency. BOFT subsumes the conventional OFT method as a particular instance, thus establishing a broader orthogonal finetuning framework. We validate BOFT through extensive experiments, adapting large-scale vision transformers, language models, and text-to-image diffusion models to a diverse range of downstream tasks across vision and language domains.
\end{abstract}

\section{Introduction}

Notable advances such as ChatGPT~\cite{brown2020language,chen2021evaluating} and Stable Diffusion~\cite{rombach2022high} have highlighted the impressive generalization capacity of contemporary large foundation models. The substantial improvement in these models’ capabilities has been matched by exponential growth in their parameter sizes

(\emph{e.g.}, GPT-3 comprises approximately 175 billion parameters). As a result, the feasibility of training such foundation models from scratch has diminished considerably for most research groups. Progress in the area now hinges on effective strategies for tailoring existing foundation models to specialized tasks, circumventing the need for complete retraining. This necessity calls for methods that efficiently adapt pre-trained models to new, downstream applications. Three principal techniques are commonly adopted for efficient adaptation: \emph{model finetuning}~\cite{hulora2022,zaken2022bitfit,guo2020parameter,raffel2020exploring,qiu2023controlling,zhang2022adaptive,chavan2023one}, in which the base architecture is preserved while selected parameters are updated; \emph{adapter tuning}~\cite{rebuffi2017learning,houlsby2019parameter,pfeiffer2020adapterfusion,lin2020exploring,he2021towards}, where extra trainable modules are incorporated into the model; and \emph{prompt tuning}~\cite{li2021prefix,lester2021power}, which introduces trainable prefix tokens to the model inputs. Among these, model finetuning stands out due to its simplicity and effectiveness, as well as its advantage of maintaining inference speed without introducing additional latency.

Model finetuning fundamentally aims to maintain similarity between the pretrained and the finetuned model, according to specific criteria, in order to retain the model’s original knowledge. Commonly, current finetuning techniques employ a reduced learning rate, a practical method that implicitly limits divergence between the pretrained parameters and their updated counterparts. Recognizing the inherent structure in weight matrices, a more theoretically grounded strategy seeks to conserve the relational properties among these matrices—specifically, the pairwise angular relationships between neurons. This perspective underpins the development of the Orthogonal Finetuning~(OFT)~\cite{qiu2023controlling} framework, where an identical orthogonal transformation is applied to all neurons within a layer, thereby rigorously maintaining inter-neuronal angles during the finetuning stage. While OFT achieves strong generalization and convergence, especially for finetuning text-to-image diffusion models~\cite{qiu2023controlling}, it introduces a significant parameter overhead due to the size of orthogonal matrices involved. To alleviate this, OFT employs a block-diagonal arrangement, cutting down on the trainable parameter count. Nonetheless, this design choice imposes certain trade-offs—the sparsity pattern of the orthogonal matrix becomes fixed, and each block’s transformation is executed independently. As a result, although this approach enhances parameter efficiency, it can also introduce less desirable biases (for example, the block-diagonal form may constrain representational flexibility and limit the matrix’s capacity to approximate general linear transformations), and crucially, the optimal design of the sparsity pattern itself remains unresolved.

A central challenge lies in constructing a dense orthogonal matrix that remains efficient in terms of parameter count. At first, this seems prohibitive, as forming a dense orthogonal matrix of dimensionality $d$ typically requires $\mathcal{O}(d^2)$ parameters. We address this by adopting an unconventional strategy: assembling a dense orthogonal matrix from a sequence of factorized sparse matrices. This is motivated by the observation that, through increased computational effort, one can effectively decrease the parameter space. Because our orthogonal matrix is the product of several sparse matrices, each training iteration necessitates multiple multiplications. Framing this matrix decomposition as an information transmission task, we view the formation of a dense orthogonal matrix as the propagation of information over a grid-based graph. This viewpoint reveals a range of approaches for sparsely factorizing matrices that constrain the parameter set while preserving the capacity to form dense matrices.

To facilitate parameter savings within this information transmission paradigm, we draw from the butterfly patterns found in the Cooley-Tukey fast Fourier transform algorithm~\cite{cooley1965algorithm}, where information efficiently flows between nodes~\cite{klasing1994broadcasting}. In the Cooley-Tukey method, the butterfly network naturally leads to a particular sparse matrix decomposition, referred to as butterfly factorization. By employing butterfly factorization, it becomes feasible to express a $d$-dimensional dense matrix as a product of $\mathcal{O}(\log d)$ sparse matrices, each containing only $\mathcal{O}(d)$ nonzero elements. Maintaining orthogonality in each sparse factor enables an efficient orthogonal parameterization comprising just $\mathcal{O}(d\log d)$ parameters, a drastic reduction relative to conventional approaches. Building on this compact orthogonal form, we introduce a new finetuning technique named Orthogonal Butterfly~(BOFT), which encompasses OFT as a specific instance and establishes a broad orthogonal finetuning framework. Both the block-diagonal and butterfly structures share a foundation of sparsity, thereby reducing the trainable parameter count within orthogonal matrices. This stands in contrast—and also in relation—to the use of low-rank architectures in LoRA~\cite{hulora2022}; both sparse and low-rank matrices represent primary classes of structured matrices~\cite{candes2011robust} that enhance parameter efficiency.

In contrast to OFT's reliance on a block-diagonal form to mediate between expressiveness and regularization, BOFT adopts the butterfly structure, facilitating a more gradual transition between matrices within the full orthogonal group (maximal expressivity) and the identity matrix (maximal regularity). This construction allows us to target a more compact hypothesis space inside the orthogonal group, better suited for subsequent applications. Given the butterfly structure’s prominence in computationally efficient linear transformations—including the discrete Fourier and cosine transforms—we contend that our structurally informed parameter-efficient strategy induces a helpful inductive bias, ultimately enhancing model generalization and reducing the risk of overfitting. This reasoning stems from the unique ability of the butterfly structure to realize a wide variety of traditional linear transforms, an ability unattainable for OFT’s block-diagonal form. Our main contributions are as follows:

\begin{itemize}[leftmargin=*,nosep]
    \item We approach the challenge of parameter-efficient orthogonal finetuning through a novel information transmission paradigm, recasting the design of dense, parameter-efficient orthogonal matrices as a problem of transmitting information across a grid-based graphical model.
    \item Motivated by the butterfly structures found in the Cooley-Tukey algorithm, we introduce Orthogonal Butterfly, a method for parameter-efficient orthogonal finetuning operating within the information transmission framework.
    \item We present several theoretical justifications demonstrating how BOFT can drastically curtail the number of learnable parameters without sacrificing expressive power or training stability. Owing to its matrix decomposition properties, BOFT additionally supports an interesting form of weight interpolation (refer to Figure~\ref{fig:boft_interpolation}).
    \item We extend orthogonal finetuning to a diverse array of domains for the first time, beyond the use case of controllable text-to-image generation~\cite{qiu2023controlling}, establishing its viability as a general-purpose model adaptation technique. In our empirical studies spanning computer vision and natural language processing tasks, BOFT surpasses prevailing state-of-the-art parameter-efficient methods by a significant margin, attesting to its superior parameter efficiency and generalization capabilities.
\end{itemize}

\textbf{Parameter-efficient finetuning (PEFT)}. With the continued scaling of foundation models (\emph{e.g.}, \cite{brown2020language,radford2021learning,rombach2022high,kirillov2023segment}), there is increasing attention on adapting these large models for specific tasks using methods that require minimal updates to model parameters~\cite{treviso2023efficient,ding2023parameter,lialin2023scaling}. Within the spectrum of PEFT approaches~\cite{houlsby2019parameter,aghajanyan2020intrinsic,hulora2022,edalati2022krona,wang2022adamix,gheini2021cross,zaken2022bitfit,guo2020parameter,sung2021training,ansell2022composable,lester2021power,li2021prefix,vu2022spot,he2021towards,mao2021unipelt,karimi2021compacter,liu2022few,sung2022lst,chen2022parameter,jia2022visual,chen2022adaptformer,zhang2022neural,jie2023fact,lian2022scaling,luo2023towards}, reparameterization-based strategies~\cite{aghajanyan2020intrinsic,hulora2022,edalati2022krona} are particularly relevant to our study. LoRA~\cite{hulora2022}, for example, adapts pretrained weight matrices by introducing the sum of the original weight and the product of two low-rank matrices, achieving strong results on a variety of NLP benchmarks. While LoRA fixes the rank of the added matrices, other methods~\cite{zhang2022adaptive,valipour2022dylora,zhang2023increlora} allocate parameter resources more flexibly by varying the rank across layers. The straightforward nature of this low-rank reparameterization has led to widespread adoption~\cite{dettmers2023qlora,zi2023delta,chavan2023one}. Drawing inspiration from research linking hyperspherical energy and model generalization~\cite{liu2018learning,liu2021orthogonal}, \cite{qiu2023controlling} proposed orthogonal finetuning, which takes a different yet effective approach to finetune text-to-image diffusion models. In this method, an orthogonal transformation is learned for neurons within the same layer, leading to improved generalization and greater training stability in comparison to LoRA. However, OFT often introduces more trainable parameters relative to LoRA, motivating the need to enhance its parameter efficiency. Additionally, the utility of OFT beyond its initial use case of text-to-image diffusion adaptation~\cite{qiu2023controlling} remains under-explored. The introduction of BOFT addresses this by employing butterfly factorization to improve OFT's parameter efficiency, thereby enabling the exploration of orthogonal finetuning across a broader set of adaptation tasks.

\textbf{Butterfly structures}. The radix-2 Cooley-Tukey algorithm offers a recursive scheme to break down an $N$-point discrete Fourier transform into two smaller transforms of size $\frac{N}{2}$, resulting in a computational structure characterized by a sequence of sparse matrix multiplications, collectively known as a butterfly matrix. Butterfly matrices have found utility in a number of contexts, such as efficiently representing orthogonal matrices in order to circumvent pivoting in Gaussian elimination~\cite{parker1995random}, enhancing training stability in recurrent neural networks~\cite{jing2017tunable}, and accelerating kernel methods~\cite{munkhoeva2018quadrature}. Research by \cite{dao2019learning,dao2019kaleidoscope} demonstrated learning fast linear transformations with butterfly parameterizations, while \cite{chen2022pixelated,dao2022monarch} integrated butterfly matrices to enable efficient neural network training. Additionally, butterfly structures have been leveraged for rapid matrix-vector multiplication~\cite{michielssen1996multilevel,o2010algorithm}, data-sparse matrix compression~\cite{li2015butterfly}, and for optimizing communication in network transmission~\cite{klasing1994broadcasting,li2003linear,rajasingh2016transmission}. Distinct from prior work, our focus lies in leveraging butterfly structures to increase the parameter efficiency specifically within the OFT framework.

\section{An Information Transmission View on Orthogonal Finetuning}

We begin by outlining the foundational concepts behind OFT. In order to adapt the pretrained weight matrix, OFT reformulates the updated weight matrix as the multiplication of a learnable orthogonal matrix with the static pretrained weight matrix. This approach contrasts with that of LoRA, which employs an additive low-rank matrix to adjust the weights. OFT instead leverages a multiplicative orthogonal matrix for weight updates. LoRA achieves efficiency in the number of parameters by adopting a low-rank design, whereas standard OFT introduces a block-diagonal orthogonal configuration~\cite{qiu2023controlling}; notably, smaller diagonal block sizes correspond to higher parameter efficiency for OFT. Figure~\ref{fig:comp_lora} provides an intuitive side-by-side comparison. The underlying rationale for utilizing orthogonal transformations in tuning the weight matrix is to retain the angular relationships between neurons~\cite{liu2018learning,liu2021orthogonal,qiu2023controlling}, which helps maintain semantic knowledge obtained during pretraining. 

Specifically, OFT seeks to learn an orthogonal matrix $\bm{R}\in\mathbb{R}^{d\times d}$ for a given pretrained linear layer $\bm{W}^0\in\mathbb{R}^{d\times n}$, altering the forward computation from $\bm{z}=(\bm{W}^0)^\top\bm{x}$ to $\bm{z}=(\bm{R}\bm{W}^0)^\top\bm{x}$, with input $\bm{x}\in\mathbb{R}^{d}$ and output $\bm{z}\in\mathbb{R}^{n}$. To ensure the weights are unchanged initially, $\bm{R}$ is set to the identity matrix at initialization. Maintaining the orthogonality of $\bm{R}$ during the fine-tuning process relies on Cayley parameterization~\cite{qiu2023controlling,liu2021orthogonal}, i.e., $\bm{R}=(\bm{I}+\bm{Q})(\bm{I}-\bm{Q})^{-1}$, where $\bm{Q}$ is a skew-symmetric matrix ($\bm{Q}=-\bm{Q}^\top$). 

To further reduce the parameter count, the block-diagonal pattern is adopted by expressing $\bm{R}$ as $\text{diag}(\bm{R}_1,\bm{R}_2,\cdots,\bm{R}_r)$, where each $\bm{R}_i\in\mathcal{R}^{b\times b}$ is a small orthogonal matrix, and $br=d$. This parameter saving is achieved at the cost of imposing an implicit constraint: the neuron's dimension (that is, each column vector in $\bm{W}^0$) is divided into $r$ separate groups, with each group being independently transformed using its respective orthogonal matrix. While empirical evidence supports the effectiveness of block-diagonal structures, this method of grouping based strictly on index lacks conceptual justification, which highlights the potential value of employing dense orthogonal matrices. This brings forth a natural inquiry: \emph{Is it possible to obtain a dense orthogonal matrix while maintaining parameter efficiency?}

In response to this issue, we suggest constructing a dense orthogonal matrix by sequentially multiplying several sparse orthogonal matrices. To achieve a cohesive yet clear framework for analyzing sparsity structures in orthogonal matrix factorizations, we reinterpret the creation of a dense orthogonal matrix in OFT through the lens of information transmission. Concretely, the generation of a dense matrix $\bm{R}\in\mathbb{R}^{d\times d}$ as the product of $m$ square matrices, written as $\thickmuskip=2mu \medmuskip=2mu \bm{R}=\bm{B}_m\bm{B}_{m-1}\cdots\bm{B}_1$, can be likened to the process of transmitting information across a grid of $d\times (m+1)$ nodes, as depicted in Figure~\ref{fig:info_trans}. The rationale for employing the information transmission analogy stems from recognizing that a $d$-dimensional dense square matrix represents a full set of connections between two layers of $d$ nodes each. If $R_{ij}$ in $\bm{R}$ equals zero, it implies the absence of a communication path from node $j$ to node $i$; conversely, a nonzero $R_{ij}$ signals that transmission is possible. Consequently, expressing $\bm{R}$ as a sequence of factors $\bm{B}_m\bm{B}_{m-1}\cdots\bm{B}_1$ can be viewed as enacting a series of information exchanges governed by the connectivity graphs induced by each $\bm{B}_i$. Information flow is initiated by $\bm{B}_1$ and culminates with $\bm{B}_n$. For illustration, Figure~\ref{fig:info_trans} presents a specific example with $\bm{R}=\bm{B}_5\bm{B}_4\bm{B}_3\bm{B}_2\bm{B}_1$, alongside their respective sparsity patterns and associated graph. The graph therein is constructed by “unrolling” the process of matrix multiplication. In this induced graph, each matrix $\bm{B}_i$ determines the connections from nodes at level $i$ to those at level $i+1$, where the $(j_1, j_2)$ entry of $\bm{B}_i$ encodes the existence (nonzero) or absence (zero) of a directed edge from the $j_2$-th node in level $i$ to the $j_1$-th node in level $i+1$. For the matrix $\bm{B}_5\bm{B}_4\bm{B}_3\bm{B}_2\bm{B}_1$ to be fully dense, it is essential that every node at the initial level can transmit to every node at the sixth level. Restricting attention to $\bm{R}=\bm{B}_4\bm{B}_3\bm{B}_2\bm{B}_1$—mapping sources at level one to receivers at level five—we see that transmission from node 1 to node 3 is blocked; this corresponds to a zero at position $(3,1)$ in the sparsity pattern. This analogy similarly holds for shorter products such as $\bm{R}=\bm{B}_3\bm{B}_2\bm{B}_1$ and $\bm{R}=\bm{B}_2\bm{B}_1$, with the induced graphs mirroring the matrices’ sparsity structures. In a general sense, ensuring that a matrix $\bm{R}\in\mathbb{R}^{d\times d}$ is dense necessitates that the sequence of $m$ matrices $\bm{B}_m,\ldots,\bm{B}_1$ collectively induces a network of directed edges on a $d\times (m+1)$ grid. Each directed edge must connect nodes between neighboring columns (levels), and the structure must permit information from any starting node at level one to be relayed to all nodes at level $m+1$.

The information transmission framework provides valuable insight into the sparsity characteristics of factorization matrices within OFT. As shown in Figure~\ref{fig:blk_diag}, the original OFT imposes a block-diagonal form on $\bm{R}$. Although this arrangement decreases the parameter count that must be learned, it inherently prevents $\bm{R}$ from attaining full density. The central objective is thus to synthesize a dense orthogonal matrix by combining $m$ sparse orthogonal matrices, while utilizing the minimal number of effective trainable parameters. According to the information transmission perspective, two general requirements are vital for achieving this: \emph{(i)} \textbf{dense connectivity}, wherein each node at the initial level maintains at least one connection to every node at the final level; and \emph{(ii)} \textbf{minimum free edges}, signifying the minimization of total edges compatible with orthogonality constraints. Orthogonality introduces nuanced requirements regarding edges linking adjacent levels. Specifically, for every matrix $\bm{B}_i$ to satisfy the full-rank condition—a prerequisite for orthogonality—it is necessary to establish $d$ edges to guarantee a bijective mapping between all nodes at the $i$-th and $(i=1)$-th levels. Consequently, the edges bridging adjacent layers must number at least $d$ (for instance, 4 as in Figure~\ref{fig:info_trans}). Since these $d$ edges are required for orthogonality and are not trainable (for example, in a $d\times d$ orthogonal matrix with $d$ non-zero elements, these must assume values of $\pm 1$), they are excluded from the count of edges that may be learned. The orthogonality property further reduces parameter requirements in comparison to unconstrained matrices, introducing extra interdependence among edges. For instance, within a $2\times 2$ orthogonal matrix, setting the $(1,1)$ position to zero necessitates a zero in the $(2,2)$ position as well; that is, omitting a single edge may entail the absence of another. As such, the orthogonality constraint can result in the addition or removal of edges depending on the configuration. The information transmission framework proves especially effective in the context of OFT by offering a visualization of the non-zero structure of the resulting orthogonal matrix; here, our interest centers solely on the locations of trainable non-zero elements in $\bm{R}$, with their precise numeric values being irrelevant.

A straightforward dense linkage between two layers requires $\mathcal{O}(d^2)$ connections (\emph{i.e.}, implementing a full $d\times d$ orthogonal matrix), resulting in $d^2-d$ connections in total (for $d=4$, this translates to 12 connections). As illustrated in Figure~\ref{fig:info_trans}, a specific matrix factorization is shown that utilizes $10$ connections overall—fewer than those needed for a fully dense orthogonal mapping. This graphical approach allows us to analyze the parameter-efficiency characteristics of OFT, and, with this viewpoint, generating feasible matrix factorizations becomes more straightforward. Motivated by the Cooley-Tukey algorithm, we draw on the structure known as butterfly graphs~\cite{cooley1965algorithm}, which achieve efficient all-to-all connectivity between $d$ sources and $d$ sinks using only $\mathcal{O}(d\log d)$ connections. For instance, whereas the example topology in Figure~\ref{fig:info_trans} consumes $10$ connections to realize dense transfer, the butterfly topology requires just $8$. The following section explains how the butterfly design increases parameter efficiency.

\section{Orthogonal Parameterization by Butterfly Factorization}

Initially devised for the Cooley-Tukey algorithm to expedite the fast Fourier transform, the butterfly structure possesses the notable property that alterations in the frequency domain can propagate globally in the spatial domain. This notion aligns well with the information transfer objective—namely, allowing every input node at the first stage to communicate with every output node at the last stage. The butterfly topology is also extensively used in high-performance computer networking~\cite{solihin2015fundamentals,li2011linear} to promote effective data exchange. Letting $k\geq 2$ be a power of $2$, we formally introduce the \emph{butterfly factor} $\bm{B}^F(k)$ as follows:

\begin{equation}\label{eq:but_factor}
\begin{aligned}
    \bm{B}^F(k)=\begin{bmatrix}
    \text{diag}(\bm{d}_1) & \text{diag}(\bm{d}_2) \\
    \text{diag}(\bm{d}_3) & \text{diag}(\bm{d}_4)
    \end{bmatrix}\in\mathbb{R}^{k\times k},
\end{aligned}
\end{equation}

with $\bm{d}_i\in\mathbb{R}^{\frac{k}{2}}$ for all $i$. For $d=2^N$, the recursive definition of the $d$-dimensional \emph{butterfly matrix} $\bm{B}(d)\in\mathbb{R}^{d\times d}$ is then

\begin{equation}
\begin{aligned}
    \!\!\bm{B}(d)=\tilde{\bm{B}}(d,d)\!\cdot\!\begin{bmatrix}
    \bm{B}_1(\frac{d}{2})\!\!\!\!\! & \bm{0} \\
    \bm{0}\!\!\!\!\! &\bm{B}_2(\frac{d}{2})
    \end{bmatrix}= \tilde{\bm{B}}(d,d)\tilde{\bm{B}}(d,\frac{d}{2})\cdots\tilde{\bm{B}}(d,2),
\end{aligned}
\end{equation}

Here, $\bm{B}_1(\frac{d}{2})$ and $\bm{B}_2(\frac{d}{2})$ denote butterfly matrices with dimensionality $\frac{d}{2}$. We introduce the \emph{butterfly component} as $\thickmuskip=2mu \medmuskip=2mu \tilde{\bm{B}}(d,k)=\text{diag}(\bm{B}^F_1(k),\cdots,\bm{B}^F_{d/k}(k))$, a block-diagonal matrix of dimension $d\times d$, where each block is of size $k$ and every $\bm{B}^F_i(k)$ (for all $i$) corresponds to the butterfly factors from Equation~\ref{eq:but_factor}. Utilizing this construction, we are able to represent an orthogonal matrix by parameterizing it with butterfly matrices. This only requires ensuring that each multiplicative block $\tilde{\bm{B}}(d,k)$ within the overall butterfly matrix $\bm{B}(d)$ is itself orthogonal. We begin by examining $\tilde{\bm{B}}(d,2)$, the block-diagonal form where each block has size $2$. Orthogonality for $\tilde{\bm{B}}(d,2)$ can be enforced using the Cayley transform (or by applying $2$-dimensional rotations) on individual blocks, as demonstrated in prior work~\cite{liu2021orthogonal,qiu2023controlling}. Since the sparsity structure of each butterfly component is essentially a permutation of the pattern in $\tilde{\bm{B}}(d,2)$, every butterfly component can be parameterized as an orthogonal matrix in an analogous fashion. As a result, this framework leads to a highly efficient orthogonal matrix parameterization using numerous $2\times 2$ orthogonal matrices as building blocks. To extend butterfly matrices following~\cite{chen2022pixelated}, we introduce the block butterfly component $\tilde{\bm{B}}^b(d,k)$, where each entry in $\bm{d}_i$ for all $i$ is upgraded to a $b\times b$ matrix. For the block butterfly component $\tilde{\bm{B}}^b(d,2)$ to remain orthogonal, each $2b\times 2b$ block must be parameterized as an orthogonal matrix. The block-sparsity pattern for subsequent butterfly components $\tilde{\bm{B}}^b(d,k)$ with $k>2$ can be regarded as a permutation of the block pattern in $\tilde{\bm{B}}^b(d,2)$, permitting orthogonality to be enforced analogously. With these elements combined, the BOFT forward computation proceeds as follows:

\begin{equation}
    \begin{aligned}\nonumber
    \bm{z}=\big(\bm{R}(m,b)\cdot\bm{W}^0\big)^\top\bm{x},~~~~\text{s.t.}~\bigg{\{}\bm{R}(m,b)=\prod_{i=1}^m \tilde{\bm{B}}^b_{(d,i)}~~\&~~\underbrace{\big(\tilde{\bm{B}}^b_{(d,j)}\big)^\top\tilde{\bm{B}}^b_{(d,j)}=\tilde{\bm{B}}^b_{(d,j)}\big(\tilde{\bm{B}}^b_{(d,j)}\big)^\top\!=\bm{I}_d}_{\forall j\in [1, m]}\bigg{\}},
    \end{aligned}
\end{equation}

Here, for notational clarity, we use $\tilde{\bm{B}}^b(d,2^{m - i+1}) = \tilde{\bm{B}}^b_{(d,i)}$, with $\bm{I}_d$ being the $d$-dimensional identity matrix. The orthogonal matrix $\bm{R}(m,b)\in\mathbb{R}^{d\times d}$ is constructed as a product of multiple orthogonal butterfly factors. To simplify references, BOFT equipped with $\bm{R}(m,\frac{b}{2})$ is denoted as BOFT($m$,$b$), where $b\geq 2$. When $\thickmuskip=2mu \medmuskip=2mu m=1$, BOFT($1$,$b$) collapses to the block-diagonal OFT~\cite{qiu2023controlling} with block size $b$. The setting BOFT($1$,$d$) becomes equivalent to the standard OFT~\cite{qiu2023controlling} that uses an unrestricted dense orthogonal matrix. For BOFT($\log \frac{2d}{b}$,$b$), the block butterfly matrix $\bm{B}^{\frac{b}{2}}(d)$ is used for $\bm{R}$, generating a dense orthogonal matrix $\bm{R}$. In a general case, the total number of trainable parameters for BOFT($m$,$b$) is $\frac{1}{2}(b-1)dm$ when finetuning a linear layer sized $d\times n$. Should the butterfly matrix be adopted, i.e., for $m=\log d, b=2$, BOFT only requires $\mathcal{O}(d\log d)$ parameters. By contrast, a classical OFT that employs a fully dense orthogonal matrix necessitates $\mathcal{O}(d^2)$ parameters, while the block-diagonal OFT with $r$ blocks takes $\mathcal{O}(bd)$. Therefore, to construct a dense orthogonal matrix, conventional OFT demands $b=d$, whereas BOFT is flexible with respect to $b$.

\textbf{Identity Initialization in BOFT}. Standard finetuning protocols begin from the original pretrained weights to preserve the behavior of the initial model. For instance, LoRA employs a zero initialization for low-rank parameter increments. Similarly, BOFT initializes every butterfly component to the identity matrix (meaning, the corresponding skew-symmetric matrices in the Cayley transform are set to zero).

\textbf{Multiplicative Dropout Strategy for BOFT}. The LoRA approach~\cite{hulora2022} leverages Dropout for the low-rank parameter updates to mitigate overfitting. While classic Dropout~\cite{srivastava2014dropout} is readily applicable to LoRA, it cannot be directly used in BOFT due to the multiplicative nature of its updates. To resolve this, we introduce a multiplicative Dropout variant for BOFT. Since $\bm{R}(m,b)$ is composed of $m$ orthogonal butterfly factors, and each can be rearranged as $2b\times 2b$ block-diagonal orthogonal matrices, our multiplicative Dropout operates by first selecting, at random, $p_1$ percent of the butterfly factors and $p_2$ percent of the diagonal blocks within each butterfly factor. The chosen blocks are then replaced with identity matrices.

\section{Intriguing Insights and Discussions}

\textbf{Expressivity of BOFT}. The butterfly structure combined with permutations is capable of exactly reconstructing numerous well-known fast linear transforms~\cite{dao2019learning,dao2019kaleidoscope} (\emph{e.g.}, fast Fourier transform, Hadamard transform). However, the extent to which our orthogonal butterfly matrix can serve as a good approximation for arbitrary orthogonal matrices has yet to be fully understood. To investigate this, we perform an empirical study by attempting to approximate a randomly generated dense orthogonal matrix~\cite{anderson1987generation} of size $512\times 512$. Figure~\ref{fig:para_eff} presents the average results across 10 different random initializations. The y-axis corresponds to the approximation error, while the x-axis shows the count of effective trainable parameters. Curves of identical color correspond to BOFT configurations with the same block size, with the leftmost marker reflecting the error achieved by BOFT($1$,$b$) (\emph{i.e.}, the vanilla block-diagonal OFT of block size $b$). BOFT consistently demonstrates superior parameter efficiency in comparison to OFT. For instance, BOFT($9$,$2$) achieves greater expressivity than BOFT($1$,$16$) despite requiring substantially fewer parameters. Generally, BOFT settings that employ smaller $b$ and higher $m$ attain better parameter efficiency; for example, BOFT($6$,$4$) matches the approximation error of BOFT($2$,$16$) using significantly fewer parameters. Overall, the butterfly matrix embodies a more constrained but structured subset of the orthogonal group relative to block-diagonal approaches, rendering BOFT provably more expressive than OFT at equal block sizes.

\begin{theorem}[Expressivity of BOFT]\label{thm:exp_boft}
    With a fixed block size, BOFT possesses strictly greater expressiveness than OFT. Given any orthogonal matrix of dimension $d$, it can be approximated by products of butterfly matrices as $\bm{B}_{d-1,1}(d)\bm{B}^\top_{d-1,2}(d)\cdots\bm{B}_{1,1}(d)\bm{B}^\top_{1,2}(d)$, where each $\bm{B}_{i,j}(d),\forall i,\forall j$ denotes a butterfly matrix.
\end{theorem}

According to Theorem~\ref{thm:exp_boft}, BOFT naturally admits an extension whereby the resulting orthogonal matrix generalizes to $\thickmuskip=2mu \medmuskip=2mu \bm{R}^G(m_1,b_1,m_2,b_2,l)=\bm{R}_{l,1}(m_1,b_1)\bm{R}_{l,2}^T(m_2,b_2)\cdots\bm{R}_{1,1}(m_1,b_1)\bm{R}_{1,2}^T(m_2,b_2)$, with each $\bm{R}_{i,j}^T(m,b)$ denoting a BOFT orthogonal block. When the parameters satisfy $m_1=m_2=\log d$, $b_1=b_2=2$, and $l=d-1$, the composition $\bm{R}^G(m_1,b_1,m_2,b_2,l)$ is able to represent any orthogonal matrix (the full orthogonal group), corresponding to the structure known as the kaleidoscope hierarchy~\cite{dao2019kaleidoscope}. Nevertheless, it is important to recognize that heightened expressivity does not automatically translate into superior finetuning performance, as even fully expressive finetuning often leads to subpar results. Achieving effective model finetuning hinges on balancing expressivity and regularization. BOFT expands the parameter space during finetuning while embedding structural constraints, supporting a more favorable compromise between expressivity and regularity.

\textbf{Spectral properties}. Compared to LoRA, orthogonal finetuning methods more effectively maintain the spectral properties of the pretrained weight matrix $\bm{W}^0$, as they leave its spectral norm unaltered. To see this, observe that for the singular value decomposition $\bm{W}^0 = \bm{U}\bm{\Sigma}\bm{V}^\top$ (where $\bm{U}$ and $\bm{V}$ are orthogonal, and $\bm{\Sigma}$ is diagonal containing the singular values), both OFT and BOFT involve applying an orthogonal matrix $\bm{R}$ on the left, resulting in updated weights of the form $\bm{R}\bm{U}\bm{\Sigma}\bm{V}^\top$. This operation preserves all singular values, and thus the spectral norm remains the same as that of the original matrix $\bm{W}^0$. Maintaining the spectral norm in this way has been found to contribute substantially to stable training dynamics and improved model generalization~\cite{miyato2018spectral,yoshida2017spectral}. For a deeper discussion of these mathematical aspects, refer to Appendix~\ref{app:sn_property}.

\textbf{Orthogonal finetuning as learning bilinear similarity}. From another perspective, BOFT can be interpreted as optimizing a bilinear similarity function of the form $\bm{w}^0_i\bm{R}\bm{x}$, where $\bm{w}^0_i$ denotes the $i$-th neuron (column) of $\bm{W}^0$. Under this framework, BOFT equates to learning the similarity matrix $\bm{R}$, but under a strict orthogonality constraint. This framing draws a clear connection to research in distance metric learning~\cite{xing2002distance} and studies on bilinear forms~\cite{roman2005advanced}.

\textbf{Inductive bias and generalization in BOFT}. Within BOFT, the matrix $\bm{R}(m,b)$ typically lies in a specific, structured subgroup of the full orthogonal group, effectively narrowing the model's hypothesis space. This imbued structure acts as an inductive bias. We posit that the particular kind of inductive bias arising from butterfly factorization serves to improve generalization: this is because its pattern reflects that of numerous well-known linear transforms—like the discrete Fourier, sine/cosine, and Hadamard transforms~\cite{dao2019kaleidoscope}. Additionally, the use of sparse factorization inherent in BOFT may introduce further, implicit forms of inductive bias~\cite{gunasekar2017implicit,li2018algorithmic,liu2019neural}.

\textbf{Relation to butterfly-based sparse training}. Numerous studies~\cite{dao2019learning,dao2019kaleidoscope,chen2022pixelated,dao2022monarch} have investigated sparse training strategies that utilize butterfly parameterization. These works generally center on the direct reparameterization of weight matrices through butterfly structures and the end-to-end training of neural networks using these parameterizations. In contrast, \cite{dao2022monarch} examines the process of finetuning pretrained weights by projecting them onto a butterfly matrix variant, followed by the optimization of these projected subspaces for specific downstream applications. BOFT introduces an alternative finetuning methodology wherein the transformation of weights leverages layer-shared matrices, marking a notable departure from earlier approaches.

\input{rebuttal-results}

\section{Concluding Remarks and Limitations}

In this paper, we introduce Orthogonal Butterfly, a universally applicable and parameter-efficient finetuning technique for foundation models, founded on the butterfly architecture. Our central innovation for enhancing parameter-efficiency lies in expressing a dense orthogonal matrix as a product of several sparse orthogonal matrices. To facilitate the construction of valid matrix decompositions, we develop a graphical information transmission perspective. Within this framework, we observe that the butterfly arrangement offers an effective solution to the challenge of sparse orthogonal matrix factorization. We empirically establish the utility of BOFT across a range of settings, including the finetuning of large language models, large vision models, and text-to-image generative models. The experimental results also underscore the generality and advantage of BOFT over existing finetuning approaches.

Nevertheless, BOFT is not without shortcomings. The design, which constructs the final orthogonal matrix as a composition of multiple orthogonal factors, introduces a slightly higher training runtime compared to OFT. Addressing this computational overhead during training is an unresolved issue. On a positive note, once finetuning is complete, the orthogonal matrices derived through BOFT can be merged into the pretrained network without incurring any additional inference time. Additionally, it remains unclear if the butterfly network represents the optimal solution for information transmission efficiency. Our proposed information transmission viewpoint paves the way for cross-disciplinary inspiration, notably from computer networking research, where network topology efficiency is an established topic. We anticipate that alternative, possibly more efficient, network topologies could serve as promising candidates for constructing dense orthogonal matrices.